% !TEX program = xelatex
\documentclass[12pt]{article}
\usepackage[margin=28mm]{geometry}
\usepackage{hyperref}
\usepackage{amsmath, amsthm, amssymb, amsfonts}
\usepackage{tikz-cd}
\usepackage{fontspec}
\usepackage{unicode-math}
\usepackage{xeCJK}
\setmainfont{Latin Modern Roman}
\setmonofont{Latin Modern Mono}
\setCJKmainfont{Noto Serif CJK JP}
\usepackage{listings}
\lstset{basicstyle=\ttfamily\footnotesize, columns=fullflexible, breaklines=true, showstringspaces=false}

\title{lean\_kernel\_image\_task — 日本語ノート}
\author{TeX generated by ChatGPT-5 Pro\\Lean source generated by CODEX (OpenAI)}
\date{2025-09-10}

\begin{document}
\maketitle

\begin{abstract}
Bourbaki の精神（順序対と射影）:
  代数的構造は射（写像）で読む。核 `ker f` は「像が 0 に射影される元の集合」、
  像 `range f` は「ある元から f によって到達する元の集合」。加法やスカラー倍の
  閉性は、射の加法性・線形性（`map_add`, `map_smul`）を各成分へ適用することで
  証明できる。
\end{abstract}

\section*{生成元の明示（Provenance)}
\begin{itemize}
  \item \textbf{TeX 生成}: ChatGPT-5 Pro
  \item \textbf{Lean ソース生成}: CODEX (OpenAI)
\end{itemize}

\section*{概説}
本ノートはアップロードされた Lean ソースから、定義・定理の先頭行を抽出して表にまとめ、
併せて一般的な可換図式（課題名に応じて合成・核と像・双対空間の評価写像など）を掲示します。
著者名（author）には「TeXは ChatGPT-5 Pro により生成」「Lean ファイルは CODEX (OpenAI) により生成」を明記しています。

\section*{サマリー}
Total lines: 86\\
Defs: 0\\
Lemmas: 0\\
Theorems: 3\\
Structures: 0\\
Inductives: 0

\section*{図式（例）}

\[
\begin{tikzcd}
A \arrow[r, "f"] \arrow[d, "g"'] & B \arrow[d, "h"] \\
C \arrow[r, "k"'] & D
\end{tikzcd}
\]



\[
\begin{tikzcd}
0 \arrow[r] & \ker f \arrow[r, hook, "i"] & V \arrow[r, "f"] & \mathrm{im}(f) \arrow[r, two heads, "p"] & 0
\end{tikzcd}
\]


\section*{定理・定義の一覧（先頭行）}
\begin{center}
\begin{tabular}{lll}
\textbf{種別} & \textbf{名前} & \textbf{シグネチャ} \\ \hline
theorem & \texttt{ker\_is\_submodule} & \texttt{{R V W : Type*} [Semiring R] [AddCommMonoid V] [AddCommMonoid W] [Module R V] [Module R W] (f : V →ₗ[R] W) : ∀ (v w : V), v ∈ LinearMap.ker f → w ∈ LinearMap.ker f → (v + w) ∈ LinearMap.ker f := by} \\
theorem & \texttt{range\_is\_submodule} & \texttt{{R V W : Type*} [Semiring R] [AddCommMonoid V] [AddCommMonoid W] [Module R V] [Module R W] (f : V →ₗ[R] W) : ∀ (y z : W), y ∈ LinearMap.range f → z ∈ LinearMap.range f → (y + z) ∈ LinearMap.range f := by} \\
theorem & \texttt{ker\_smul\_closed} & \texttt{{R V W : Type*} [Semiring R] [AddCommMonoid V] [AddCommMonoid W] [Module R V] [Module R W] (f : V →ₗ[R] W) (a : R) (v : V) : v ∈ LinearMap.ker f → (a • v) ∈ LinearMap.ker f := by} \\
\end{tabular}
\end{center}

\section*{Lean ソース（全文）}
\begin{lstlisting}
import Mathlib.Algebra.Module.LinearMap.Defs
import Mathlib.Algebra.Module.LinearMap.Basic
import Mathlib.Algebra.Module.Submodule.Ker
import Mathlib.Algebra.Module.Submodule.Range

/-!
  Bourbaki の精神（順序対と射影）:
  代数的構造は射（写像）で読む。核 `ker f` は「像が 0 に射影される元の集合」、
  像 `range f` は「ある元から f によって到達する元の集合」。加法やスカラー倍の
  閉性は、射の加法性・線形性（`map_add`, `map_smul`）を各成分へ適用することで
  証明できる。
-/

noncomputable section

namespace MyProjects.LinearAlgebra.A

-- 線形写像の核が部分空間であることの証明
theorem ker_is_submodule {R V W : Type*}
  [Semiring R] [AddCommMonoid V] [AddCommMonoid W] 
  [Module R V] [Module R W]
  (f : V →ₗ[R] W) :
  ∀ (v w : V), v ∈ LinearMap.ker f → w ∈ LinearMap.ker f → (v + w) ∈ LinearMap.ker f := by
  intro v w hv hw
  -- `v ∈ ker f` は `f v = 0` と同値
  have hv0 : f v = 0 := by simpa [LinearMap.mem_ker] using hv
  have hw0 : f w = 0 := by simpa [LinearMap.mem_ker] using hw
  -- 線形性（加法）で `f (v + w) = f v + f w = 0`
  have : f (v + w) = 0 := by simpa [hv0, hw0] using (LinearMap.map_add f v w)
  simpa [LinearMap.mem_ker] using this

-- 線形写像の像が部分空間であることの証明  
theorem range_is_submodule {R V W : Type*}
  [Semiring R] [AddCommMonoid V] [AddCommMonoid W]
  [Module R V] [Module R W]
  (f : V →ₗ[R] W) :
  ∀ (y z : W), y ∈ LinearMap.range f → z ∈ LinearMap.range f → (y + z) ∈ LinearMap.range f := by
  intro y z hy hz
  -- `y ∈ range f` は `∃ v, f v = y` と同値
  rcases hy with ⟨v, hv⟩
  rcases hz with ⟨w, hw⟩
  refine ⟨v + w, ?_⟩
  simpa [hv, hw] using (LinearMap.map_add f v w)

-- 発展：核の元のスカラー倍も核に属する
theorem ker_smul_closed {R V W : Type*}
  [Semiring R] [AddCommMonoid V] [AddCommMonoid W]
  [Module R V] [Module R W]
  (f : V →ₗ[R] W) (a : R) (v : V) :
  v ∈ LinearMap.ker f → (a • v) ∈ LinearMap.ker f := by
  intro hv
  -- `ker f` は部分加群なのでスカラー倍で閉じている
  simpa using ((LinearMap.ker f).smul_mem a hv)

/-! ## Bourbaki スタイルの軽い example

核：`f v = 0` と `v ∈ ker f` を往復し、`map_add` により閉性を成分的に確認。
像：`y ∈ range f` を実在元の存在で読み（射影的に）加法閉性を示す。
-/

section Examples
  variable {R V W : Type*}
  variable [Semiring R] [AddCommMonoid V] [AddCommMonoid W]
  variable [Module R V] [Module R W]

  example (f : V →ₗ[R] W) (v w : V)
      (hv : v ∈ LinearMap.ker f) (hw : w ∈ LinearMap.ker f) : (v + w) ∈ LinearMap.ker f := by
    have hv0 : f v = 0 := by simpa [LinearMap.mem_ker] using hv
    have hw0 : f w = 0 := by simpa [LinearMap.mem_ker] using hw
    have : f (v + w) = 0 := by simpa [hv0, hw0] using (LinearMap.map_add f v w)
    simpa [LinearMap.mem_ker] using this

  example (f : V →ₗ[R] W) (y z : W)
      (hy : y ∈ LinearMap.range f) (hz : z ∈ LinearMap.range f) :
      (y + z) ∈ LinearMap.range f := by
    rcases hy with ⟨v, hv⟩
    rcases hz with ⟨w, hw⟩
    refine ⟨v + w, ?_⟩
    simpa [hv, hw] using (LinearMap.map_add f v w)

  example (f : V →ₗ[R] W) (a : R) (v : V)
      (hv : v ∈ LinearMap.ker f) : (a • v) ∈ LinearMap.ker f := by
    simpa using ((LinearMap.ker f).smul_mem a hv)
end Examples

end MyProjects.LinearAlgebra.A

\end{lstlisting}

\end{document}
